\section{Introduction}\label{sec:intro}

Agents that carry out machine-learning research on their own are usually built as harnesses. The harness fixes the
order of work (plan, explore, train, predict, evaluate), checks each artefact before moving on, and keeps the best
model. The language model fills in each step but does not decide what happens next. This structure makes runs
predictable and checkable, but it is costly to build, specific to one kind of task, and spends model calls on every
step whether the model needs them or not.

Agent skills offer a lighter alternative: a short document describing how to do a task, loaded by a general-purpose agent when the task calls for it. Nothing forces the agent to follow it.

The question is whether the control a harness provides is still needed as models improve. One plausible answer is that
it depends on the model: mid-range models need the direction, while top-of-line models do as well, or better, with
guidance alone. This report tests that hypothesis, grouping models by each provider's own line-up rather than by
their results.

We hold the task, the data, the procedure and the budget of model rounds fixed, and vary three things: whether the procedure is
enforced by a harness or written as a skill, which model runs it, and the model's reasoning effort. The contributions
are:
\begin{itemize}
\item a controlled comparison of a state-machine harness and a skill on the same research task, across nine models
  and two reasoning settings (\cref{sec:results});
\item an iteration-level analysis, re-scoring the model each run would have submitted after one, two and three
  iterations, including a single-shot comparison of the two methods (\cref{sec:iterations});
\item an audit of how closely the agents followed the skill's written procedure (\cref{sec:limitations});
\item an account of cost and training time per run (\cref{sec:cost});
\item a check against the published Agentomics-ML results for GPT-4.1, the one model the two studies share
  (\cref{sec:agentomics}).
\end{itemize}
